\section{Sugerencias de trabajos futuros}

A partir de las limitaciones empíricas identificadas en este estudio, se proponen cuatro líneas de continuación: (1) \textbf{Integración de variables geoespaciales y registros administrativos}: enriquecer el vector $\mathbb{R}^{28}$ con indicadores censales a nivel de manzana e imágenes satelitales (densidad constructiva, pendiente topográfica y acceso a redes de transporte y salud en periferia), reduciendo la dependencia exclusiva del autorreporte en entrevista; (2) \textbf{Protocolo específico para hogares secundarios allegados}: dado que los núcleos sin ficha física de vivienda (\texttt{sin\_datos\_vivienda = 1}) exhiben casi el doble de pobreza ($32.97\%$) y menor discriminación predictiva, los formularios operativos del SISFOH deben vincular obligatoriamente las características físicas del hogar principal de la vivienda o emplear un submodelo específico entrenado sobre variables demográficas y laborales; (3) \textbf{Recalibración adaptativa del umbral operativo $\tau$}: implementar un esquema de actualización anual ligera del corte de decisión $\tau$ sobre submuestras rotativas para absorber choques inflacionarios en la canasta básica de Lima Metropolitana sin necesidad de reentrenar la arquitectura completa; y (4) \textbf{Despliegue como microservicio contenedorizado}: acoplar el paquete serializado \texttt{models/winner\_bundle.joblib} y el motor explicativo \texttt{src/inference.py} en una API REST sobre \textit{FastAPI} y \textit{Docker} que entregue en tiempo real el dictamen de elegibilidad y los factores locales TreeSHAP a los empadronadores municipales.
